\PassOptionsToPackage{unicode=true}{hyperref} % options for packages loaded elsewhere
\PassOptionsToPackage{hyphens}{url}
%
\documentclass[twocolumn]{article}
\usepackage{lmodern}
\usepackage{amssymb,amsmath}
\usepackage{ifxetex,ifluatex}
\usepackage{fixltx2e} % provides \textsubscript
\ifnum 0\ifxetex 1\fi\ifluatex 1\fi=0 % if pdftex
  \usepackage[T1]{fontenc}
  \usepackage[utf8]{inputenc}
  \usepackage{textcomp} % provides euro and other symbols
\else % if luatex or xelatex
  \usepackage{unicode-math}
  \defaultfontfeatures{Ligatures=TeX,Scale=MatchLowercase}
\fi
% use upquote if available, for straight quotes in verbatim environments
\IfFileExists{upquote.sty}{\usepackage{upquote}}{}
% use microtype if available
\IfFileExists{microtype.sty}{%
\usepackage[]{microtype}
\UseMicrotypeSet[protrusion]{basicmath} % disable protrusion for tt fonts
}{}
\IfFileExists{parskip.sty}{%
\usepackage{parskip}
}{% else
\setlength{\parindent}{0pt}
\setlength{\parskip}{6pt plus 2pt minus 1pt}
}
\usepackage{hyperref}
\hypersetup{
            pdfborder={0 0 0},
            breaklinks=true}
\urlstyle{same}  % don't use monospace font for urls
\usepackage[left=0.5cm,right=0.5cm,top=1.5cm,bottom=1.2cm]{geometry}
\usepackage{graphicx,grffile}
\makeatletter
\def\maxwidth{\ifdim\Gin@nat@width>\linewidth\linewidth\else\Gin@nat@width\fi}
\def\maxheight{\ifdim\Gin@nat@height>\textheight\textheight\else\Gin@nat@height\fi}
\makeatother
% Scale images if necessary, so that they will not overflow the page
% margins by default, and it is still possible to overwrite the defaults
% using explicit options in \includegraphics[width, height, ...]{}
\setkeys{Gin}{width=\maxwidth,height=\maxheight,keepaspectratio}
\setlength{\emergencystretch}{3em}  % prevent overfull lines
\providecommand{\tightlist}{%
  \setlength{\itemsep}{0pt}\setlength{\parskip}{0pt}}
\setcounter{secnumdepth}{0}
% Redefines (sub)paragraphs to behave more like sections
\ifx\paragraph\undefined\else
\let\oldparagraph\paragraph
\renewcommand{\paragraph}[1]{\oldparagraph{#1}\mbox{}}
\fi
\ifx\subparagraph\undefined\else
\let\oldsubparagraph\subparagraph
\renewcommand{\subparagraph}[1]{\oldsubparagraph{#1}\mbox{}}
\fi

% set default figure placement to htbp
\makeatletter
\def\fps@figure{htbp}
\makeatother

\usepackage{fancyhdr} \usepackage{booktabs,xcolor} \pagestyle{fancy} \renewcommand{\headrulewidth}{0pt} \rhead{2019} \cfoot{\includegraphics[width=20cm]{footer.png}} \fancypagestyle{plain}{\pagestyle{fancy}} \setlength{\headheight}{77.3pt} \setlength{\footskip}{36.5pt} \setlength{\textheight}{0.90\textheight} \pagenumbering{gobble} \usepackage[fontsize=9pt]{scrextend} \usepackage{float} \restylefloat{table} \usepackage{xcolor} \usepackage{multicol} \usepackage{array} \usepackage{colortbl} \usepackage{multirow} \usepackage{collcell} \usepackage{setspace} \usepackage{arydshln} \usepackage{caption} \captionsetup{skip=0pt} \setlength{\columnsep}{1.5cm}

\author{}
\date{\vspace{-2.5em}}

\begin{document}

\newcommand{\greynote}[1]{
    {\scriptsize
    \textcolor{darkgray}{\textit{Notes:} #1}
  }
}

\newcommand{\greysource}[1]{
    {\scriptsize
    \textcolor{darkgray}{\textit{Source:} #1}
  }
}

\newcommand{\greydisclaimer}[1]{
    {\scriptsize
    \textcolor{darkgray}{\textit{Disclaimer:} #1}
  }
}

\newcommand*{\tabindent}{\hspace{1mm}}

\hypertarget{introduction}{%
\subsubsection{\texorpdfstring{\textbf{INTRODUCTION}}{INTRODUCTION}}\label{introduction}}

\textbf{The Global Education Policy Dashboard (GEPD): An innovative tool
to measure drivers of learning outcomes in basic education}\\
GEPD uses 3 data collection instruments to report on nearly 40
indicators that operationalize the World Development Report 2018
framework to track 3 areas for progress in education- Practices,
Policies, and Politics. Using these indicators, the dashboard highlights
areas where countries need to act to improve learning outcomes and
allows a way for governments to track progress as they act to close gaps
in these areas. For more information on GEPD, please visit
\textbf{\href{https://www.worldbank.org/en/topic/education/brief/global-education-policy-dashboard}{www.worldbank.org/global-education-policy-dashboard}}

\hypertarget{figure-1.-gepd-framework-practices-policies-and-politics-expanding-on-wdr-2018-framework}{%
\paragraph{Figure 1. GEPD framework (practices, policies and politics),
expanding on WDR 2018
framework}\label{figure-1.-gepd-framework-practices-policies-and-politics-expanding-on-wdr-2018-framework}}

\begin{center}\includegraphics[width=0.8\linewidth]{/Users/kanikaverma/Documents/GitHub/gepd/Country_Reports/Data/full_circle} \end{center}

\hypertarget{instruments-of-data}{%
\subsubsection{\texorpdfstring{\textbf{INSTRUMENTS OF
DATA}}{INSTRUMENTS OF DATA}}\label{instruments-of-data}}

The \textbf{School Survey} consists of 8 modules to collect data across
200-300 schools on practices (the quality of service delivery in
schools) and de facto policy indicators. It consists of streamlined
versions of existing instruments together with new questions to fill
gaps in those instruments.\\
The \textbf{Policy Survey} collects information via interviews with
\textasciitilde{}200 officials per country at federal and regional level
to feed into the policy de jure indicators and identify key elements of
the policy framework.\\
The \textbf{Survey of public officials} collects information about the
capacity and orientation of the bureaucracy and political factors
affecting education outcomes. This survey is an education-focused
version of the civil-servant surveys from the Bureaucracy Lab, WBG.

\hypertarget{key-takeaways-2019}{%
\subsubsection{\texorpdfstring{\textbf{KEY TAKEAWAYS, \uppercase{Peru},
2019}}{KEY TAKEAWAYS, , 2019}}\label{key-takeaways-2019}}

\begin{itemize}
\tightlist
\item
  52\% learning poverty is observed at grade 4 level
\item
  Overall Grade 4 student proficiency (\textgreater{}80\% student
  knowledge) is only 33\%, attributed to low numeracy proficiency 23\%).
\item
  Low teacher content knowledge attributed to poor teacher support and
  lack of strong instructional leadership (especially principal feedback
  on classroom observation)
\item
  Low capacity for learning in Grade 1 (particularly socio-emotional
  knowledge) is attributed to low enrolment in early childhood programs
  (13\%)
\item
  A \textasciitilde{}0.8 point difference exists in the policy de-jure
  and de-facto indicators on a 5 point scale for teaching and school
  management. Major gaps are seen in teacher support, teacher motivation
  and evaluation of school management
\item
  Bureaucracy scores high on mandates and accountability but low on
  impartial decision making
\end{itemize}

\setlength\dashlinedash{0.2pt}
\setlength\dashlinegap{1.5pt}
\setlength\arrayrulewidth{0.3pt}

\hypertarget{table-1.-key-gepd-outcome-indicators}{%
\paragraph{Table 1. Key GEPD Outcome
Indicators}\label{table-1.-key-gepd-outcome-indicators}}

\begin{table}[H]
\resizebox{\columnwidth}{!}{\begin{tabular}{p{0.4\textwidth}c}
\textbf{Indicator} & \textbf{Value}\\ \hline
Learning poverty & {\cellcolor{red!15}56\%} \\\cdashline{1-2}
Proficiency by End of Primary & {\cellcolor{red!15}54\%} \\\cdashline{1-2}
Proficiency on GEPD Assessment & {\cellcolor{red!15}33\%} \\\cdashline{1-2}
\hspace{1mm}\emph{Literacy proficiency} & {\cellcolor{yellow!15}61\%} \\\cdashline{1-2}
\hspace{1mm}\emph{Numeracy proficiency} & {\cellcolor{red!15}23\%} \\\cdashline{1-2}
Proficiency by Grade 2/3 & {\cellcolor{red!15}-999\%} \\\cdashline{1-2}
Net Adjusted Enrollment Rate & {\cellcolor{red!15}-999\%} \\\hline
\end{tabular}}
\\
\setstretch{0.8}\color{darkgray}\scriptsize{\textit{Source:} UIS, GLAD, GEPD, World Bank, Peru, 2019. For information on indicators, please consult the World Bank \href{https://github.com/worldbank/GEPD}{\underline{GEPD}}, \href{https://github.com/worldbank/GLAD}{\underline{GLAD}} and \href{https://github.com/worldbank/LearningPoverty}{\underline{Learning Poverty}} repositories.}\\
\setstretch{0.8}\color{darkgray}\scriptsize{\textit{Notes:} All indicators are on a scale of 0-5 unless measured in \%. Green indicates indicator 'on-target', yellow indicates 'requires caution', red indicates 'needs improvement'}
\end{table}
\raggedbottom

\hypertarget{learning-outcomes-56-learning-poverty-33-gepd-proficiency-in-grade-4}{%
\subsubsection{\texorpdfstring{\textbf{LEARNING OUTCOMES: 56\% LEARNING
POVERTY, 33\% GEPD PROFICIENCY IN GRADE
4}}{LEARNING OUTCOMES: 56\% LEARNING POVERTY, 33\% GEPD PROFICIENCY IN GRADE 4}}\label{learning-outcomes-56-learning-poverty-33-gepd-proficiency-in-grade-4}}

Learning poverty in Peru is 5\% points worse than the average for Latin
America \& Caribbean (excluding high income) region and 27\% points
worse than the average for Upper middle income countries.

\hypertarget{figure-2.-leaning-poverty-comparison-peru}{%
\paragraph{Figure 2. Leaning poverty comparison,
Peru}\label{figure-2.-leaning-poverty-comparison-peru}}

\includegraphics[width=1\linewidth]{/Users/kanikaverma/Documents/GitHub/gepd/Country_Reports/LP figures/lp_figure_PER}

{\scriptsize
    \textcolor{darkgray}{\textit{Notes:} Grey circles represent other countries. Large yellow circle represents Peru. Other circles represent average learning poverty in Peru's region and income group.}
  }

\textbf{GEPD grade 4 assessment proficiency of 33\%} means 33\% students
score greater than 80/100 on student knowledge. Proficiency is 38 points
higher in literacy compared to numeracy, 2 points higher for boys
compared to girls, and 23 points higher in urban areas compared to rural
areas.

\hypertarget{figure-3.-gepd-grade-4-proficiency-peru}{%
\paragraph{Figure 3. GEPD Grade 4 proficiency,
Peru}\label{figure-3.-gepd-grade-4-proficiency-peru}}

\includegraphics{GEPD_Brief_2page_files/figure-latex/learning_figure-1.pdf}

\hypertarget{comparing-de-facto-practices-and-policy-levers}{%
\subsubsection{\texorpdfstring{\textbf{COMPARING DE-FACTO PRACTICES AND
POLICY
LEVERS}}{COMPARING DE-FACTO PRACTICES AND POLICY LEVERS}}\label{comparing-de-facto-practices-and-policy-levers}}

Practice indicators measure quality of service delivery in schools such
as teacher and student attendance, teacher knowledge, principal
management skills, etc. Policy lever indicators measure how well school,
personnel and student policies governing these practices are
implemented. For instance, teacher content knowledge(practice) is
impacted by implementation of teacher support policies, and monitoring
and accountability systems. Comparing de-facto practice and policy lever
indicators allows identification of low-scoring policy levers that
affect observed practice indicators.

\hypertarget{teacher-effectiveness}{%
\paragraph{\texorpdfstring{\textbf{Teacher
effectiveness}}{Teacher effectiveness}}\label{teacher-effectiveness}}

Teacher content knowledge (39\%) needs improvement. Teacher proficiency
in language (41\%) is 5 points higher than mathematics proficiency
(36\%). Teacher pedagogical skills score (34\%) needs improvement, and
teacher attendance (92\%) is on target. Teaching - Monitoring \&
Accountability policy lever scores the lowest(2.8/5).

\begin{table}[H]
\resizebox{\columnwidth}{!}{\begin{tabular}{m{3.8cm}cm{4.2cm}c}
\multicolumn{2}{c}{\textbf{Practice Indicators}} & \multicolumn{2}{c}{\textbf{Policy levers (Teaching)}}\\\hline
Content knowledge                   & {\cellcolor{red!15}39\%}   & Attraction                                & \cellcolor{yellow!15}3.6 \\\cdashline{1-4} 
\hspace{1mm}\emph{Maths proficiency}               & {\cellcolor{red!15}36\%} &                                           & \cellcolor{yellow!15}\\\cdashline{1-2}   
\hspace{1mm}\emph{Language proficiency}            & {\cellcolor{red!15}41\%} & \multirow{-2}{*}{Selection \& deployment} & \multirow{-2}{*}{\cellcolor{yellow!15}3.6} \\\cdashline{1-4}        
Pedagogical skills                               & {\cellcolor{red!15}34\%}   &                                           & \cellcolor{yellow!15}\\\cdashline{1-2}
\hspace{1mm}\emph{\% Classroom culture}            & {\cellcolor{green!15}91\%} & \multirow{-2}{*}{Support}                 & \multirow{-2}{*}{\cellcolor{yellow!15}3.1} \\\cdashline{1-4}
\hspace{1mm}\emph{\% Instruction practices}        & {\cellcolor{red!15}39\%} &                                           & \cellcolor{green!15}\\\cdashline{1-2}
\hspace{1mm}\emph{\% Socio-emotional skills}       & {\cellcolor{red!15}24\%} & \multirow{-2}{*}{Evaluation}              & \multirow{-2}{*}{\cellcolor{green!15}4} \\\cdashline{1-4}
& \cellcolor{green!15} & Monitoring \& Accountability  & \cellcolor{red!15}2.8 \\\cdashline{3-4}
\multirow{-2}{*}{Teacher Attendance}             & \multirow{-2}{*}{\cellcolor{green!15}92\%} & Intrinsic motivation    & \cellcolor{green!15}4.1 \\\hline
\end{tabular}}
\\
\setstretch{0.8}\color{darkgray}\scriptsize{\textit{Notes:} Content knowledge(\& sub-indicators) indicate \% teachers with knowledge\textgreater{80\%}.Pedagogical skills(\& sub-indicators) indicate \% teachers with proficiency 3/5 or above.}
\end{table}

\hypertarget{capacity-for-learning-in-grade-1}{%
\paragraph{\texorpdfstring{\textbf{Capacity for learning in Grade
1}}{Capacity for learning in Grade 1}}\label{capacity-for-learning-in-grade-1}}

Student proficiency in Grade 1 (53\%) needs improvement. Executive
funcion score(65)\% is the lowest knowledge sub-score. Student
attendance(93\%) is on target. Center-Based Care policy lever scores the
lowest(2.6/5).

\begin{table}[H]
\resizebox{\columnwidth}{!}{\begin{tabular}{m{3.8cm}cm{4.2cm}c}
\multicolumn{2}{c}{\textbf{Practice Indicators}} & \multicolumn{2}{c}{\textbf{Policy levers (Learners)}}\\\hline
Capacity for learning                  & {\cellcolor{red!15}53\%} & & \cellcolor{green!15}\\\cdashline{1-2}
\hspace{1mm}\emph{Numeracy score}        & \cellcolor{green!15}90   & \multirow{-2}{*}{Nutrition Programs} & \multirow{-2}{*}{\cellcolor{green!15}4.3}  \\\cdashline{1-4}  
\hspace{1mm}\emph{Literacy score}        & \cellcolor{yellow!15}77   & Health Programs                     & \cellcolor{green!15}4.2 \\\cdashline{1-4}  
\hspace{1mm}\emph{Executive score}       & \cellcolor{yellow!15}65   & Center based care                   & \cellcolor{red!15}2.6 \\\cdashline{1-4}  
\hspace{1mm}\emph{Socio-emotional score} & \cellcolor{yellow!15}75   & Caregiver Skills Capacity           & \cellcolor{red!15}2.9 \\\cdashline{1-4}  
Student Attendance                     & {\cellcolor{green!15}93\%}     & Caregiver Financial Capacity        & \cellcolor{green!15}4.5 \\\hline
\end{tabular}}
\\
\setstretch{0.8}\color{darkgray}\scriptsize{\textit{Notes:} Capacity for learning indicates \% students with knowledge\textgreater{80\%}. Sub-indicator scores refer to average subject knowledge on a 0-100 scale.}
\end{table}
\vfill\null

\hypertarget{inputs-infrastructure}{%
\paragraph{\texorpdfstring{\textbf{Inputs \&
Infrastructure}}{Inputs \& Infrastructure}}\label{inputs-infrastructure}}

Basic inputs (4.1/5) are on target. Percent of schools with access to
EdTech(78)\% is the lowest score. Basic infrastructure (3.7/5) requires
caution. Percent of schools with access to internet(54)\% is the lowest
score. Inputs \& Infrastructure - Monitoring policy lever scores the
lowest(3.7/5).

\begin{table}[H]
\resizebox{\columnwidth}{!}{\begin{tabular}{m{3.8cm}cm{4.2cm}c}
\multicolumn{2}{c}{\textbf{Practice Indicators}} & \multicolumn{2}{c}{\textbf{Policy levers(Inputs)}}\\\hline
Basic inputs & \cellcolor{green!15}4.1 & & \cellcolor{yellow!15} \\\cdashline{1-2}
\hspace{1mm}\emph{\% Blackboard}    & {\cellcolor{green!15}82\%} & & \cellcolor{yellow!15}\\\cdashline{1-2}
\hspace{1mm}\emph{\% Stationery}    & {\cellcolor{green!15}96\%} & & \cellcolor{yellow!15}\\\cdashline{1-2}
\hspace{1mm}\emph{\% Furniture}     & {\cellcolor{green!15}84\%} & & \cellcolor{yellow!15}\\\cdashline{1-2}
\hspace{1mm}\emph{\% EdTech access} & {\cellcolor{yellow!15}78\%} & \multirow{-5}{4.2cm}{Inputs and infrastructure standards} & \multirow{-5}{*}{\cellcolor{yellow!15}3.7}\\\hline
Basic infrastructure & \cellcolor{yellow!15}3.7 & & \cellcolor{yellow!15} \\\cdashline{1-2}
\hspace{1mm}\emph{\% Drinking water}    & {\cellcolor{green!15}91\%} & & \cellcolor{yellow!15}\\\cdashline{1-2}
\hspace{1mm}\emph{\% Functional toilet} & {\cellcolor{yellow!15}68\%} & & \cellcolor{yellow!15}\\\cdashline{1-2}
\hspace{1mm}\emph{\% Internet}          & {\cellcolor{green!15}96\%} & & \cellcolor{yellow!15}\\\cdashline{1-2}
\hspace{1mm}\emph{\% Electricity}       & {\cellcolor{red!15}54\%} & & \cellcolor{yellow!15}\\\cdashline{1-2}
\hspace{1mm}\emph{\% Disability access} & {\cellcolor{yellow!15}65\%} & \multirow{-6}{4.2cm}{Inputs and infrastructure monitoring} & \multirow{-6}{*}{\cellcolor{yellow!15}3.7}\\\hline
\end{tabular}}
\\
{\scriptsize
    \textcolor{darkgray}{\textit{Notes:} \% refers to \% schools with the given sub-component}
  }

\end{table}

\hypertarget{school-management-by-principals}{%
\paragraph{\texorpdfstring{\textbf{School Management by
principals}}{School Management by principals}}\label{school-management-by-principals}}

In school management, the lowest score is for principal's Instructional
Leadership(3.3/5), whereas the highest score is obtained for Operational
Management(4.3/5). School Management- Support policy lever scores the
lowest(3.5/5).

\begin{table}[H]
\resizebox{\columnwidth}{!}{\begin{tabular}{m{4cm}cm{4cm}c}
\multicolumn{2}{c}{\textbf{Practice Indicators}} & \multicolumn{2}{c}{\textbf{Policy levers(Management)}}\\\hline
Operational management                       & \cellcolor{green!15}4.3                   &                                       & \cellcolor{green!15}\\\cdashline{1-2}
Instructional Leadership                     & \cellcolor{yellow!15}3.3                   & \multirow{-2}{*}{Clarity of functions} & \multirow{-2}{*}{\cellcolor{green!15}5} \\\hline
                                             & \cellcolor{yellow!15}                    & Attraction                            & \cellcolor{green!15}4.4\\\cdashline{3-4}
\multirow{-2}{*}{Principal School knowledge}  & \multirow{-2}{*}{\cellcolor{yellow!15}3.6} & Selection \& Deployment               & \cellcolor{green!15}4.4\\\cdashline{3-4}
                                              & \cellcolor{green!15}                    & Support                               & \cellcolor{yellow!15}3.5\\\cdashline{3-4}
\multirow{-2}{*}{Principal Management skills}  & \multirow{-2}{*}{\cellcolor{green!15}4.3} & Evaluation                            & \cellcolor{green!15}4.4\\\hline
\end{tabular}}
\\
{\scriptsize
    \textcolor{darkgray}{\textit{Notes:} All scores are on a (0-5) scale.}
  }

\end{table}

\hypertarget{politics-bureaucratic-capacity-indicators}{%
\subsubsection{\texorpdfstring{\textbf{POLITICS \& BUREAUCRATIC CAPACITY
INDICATORS}}{POLITICS \& BUREAUCRATIC CAPACITY INDICATORS}}\label{politics-bureaucratic-capacity-indicators}}

Politics and bureaucratic capacity indicators measure the capacity and
orientation of the bureaucracy, as well as political factors affecting
education outcomes. The highest score in politics and bureaucratic
capacity is noted for Financing (5/5), and the lowest score is noted for
National Learning Goals (3/5).

\begin{table}[H]
\resizebox{\columnwidth}{!}{\begin{tabular}{p{0.4\textwidth}c}
\textbf{Indicator} & \textbf{Value}\\ \hline
Quality of Bureaucracy & \cellcolor{green!15}4\\\cdashline{1-2}
Impartial Decision-Making & \cellcolor{yellow!15}3.3\\\cdashline{1-2}
Mandates \& Accountability & \cellcolor{green!15}4.2\\\cdashline{1-2}
National Learning Goals & \cellcolor{yellow!15}3\\\cdashline{1-2}
Financing & \cellcolor{green!15}5\\\hline
\end{tabular}}
\\
{\scriptsize
    \textcolor{darkgray}{\textit{Notes:} All scores are on a (0-5) scale.}
  }

\end{table}
\raggedbottom

{\scriptsize
    \textcolor{darkgray}{\textit{Disclaimer:} GEPD numbers presented in this brief are based on multiple sources including GEPD instruments, UIS, GLAD and Learning Poverty indicators. For that reason, the numbers discussed here may be different from official statistics reported by governments and national offices of statistics. Such differences are due to the different purposes of the statistics, which can be for global comparison or to meet national definitions.}
  }

\end{document}
